\documentclass[12pt]{exam}
\usepackage[utf8]{inputenc}

\usepackage{amsmath,amstext,amsthm,amssymb,amsxtra}
\usepackage[top=1.5in, bottom=1.5in, left=1.25in, right=1.25in]{geometry}
\usepackage{txfonts}
\usepackage[T1]{fontenc}
\usepackage{lmodern}
\renewcommand*\familydefault{\sfdefault}
\usepackage{euler}
\usepackage{pdfsync}
\usepackage{multicol}
\newcommand{\ci}[1]{_{ {}_{\scriptstyle #1}}}
\usepackage{graphicx}
\graphicspath{ {images/} }

\newcommand{\norm}[1]{\ensuremath{\left\|#1\right\|}}
\newcommand{\abs}[1]{\ensuremath{\left\vert#1\right\vert}}
\newcommand{\ip}[2]{\ensuremath{\left\langle#1,#2\right\rangle}}
\newcommand{\p}{\ensuremath{\partial}}
\newcommand{\pr}{\mathcal{P}}

\newcommand{\pbar}{\ensuremath{\bar{\partial}}}
\newcommand{\db}{\overline\partial}
\newcommand{\D}{\mathbb{D}}
\newcommand{\B}{\mathbb{B}}
\newcommand{\Sp}{\mathbb{S}}
\newcommand{\T}{\mathbb{T}}
\newcommand{\R}{\mathbb{R}}
\newcommand{\Z}{\mathbb{Z}}
\newcommand{\C}{\mathbb{C}}
\newcommand{\N}{\mathbb{N}}
\newcommand{\Q}{\mathbb{Q}}
\newcommand{\mQ}{\mathcal{Q}}
\newcommand{\mS}{\mathcal{S}}
\newcommand{\scrH}{\mathcal{H}}
\newcommand{\scrL}{\mathcal{L}}
\newcommand{\td}{\widetilde\Delta}
\newcommand{\pw}{\text{PW}}
\newcommand{\esup}{\text{ess.sup}}
\newcommand{\Tn}{\mathcal{T}_n}
\newcommand{\Bn}{\mathbb{B}_n}
\newcommand{\rt}{\mathcal{O}}
\newcommand{\avg}[1]{\langle #1 \rangle}
\newcommand{\one}{\mathbbm{1}}
\newcommand{\eps}{\varepsilon}
\newcommand{\grad}{\nabla}

\newcommand{\La}{\langle }
\newcommand{\Ra}{\rangle }
\newcommand{\rk}{\operatorname{rk}}
\newcommand{\card}{\operatorname{card}}
\newcommand{\ran}{\operatorname{Ran}}
\newcommand{\osc}{\operatorname{OSC}}
\newcommand{\im}{\operatorname{Im}}
\newcommand{\re}{\operatorname{Re}}
\newcommand{\tr}{\operatorname{tr}}
\newcommand{\vf}{\varphi}
\newcommand{\f}[2]{\ensuremath{\frac{#1}{#2}}}

\newcommand{\kzp}{k_z^{(p,\alpha)}}
\newcommand{\klp}{k_{\lambda_i}^{(p,\alpha)}}
\newcommand{\TTp}{\mathcal{T}_p}
\newcommand{\m}[1]{\mathcal{#1}}
\newcommand{\md}{\mathcal{D}}
\newcommand{\qan}{\abs{Q}^{\alpha/n}}
\newcommand{\sbump}[2]{[[ #1,#2 ]]}
\newcommand{\mbump}[2]{\lceil #1,#2 \rceil}
\newcommand{\cbump}[2]{\lfloor #1,#2 \rfloor}

\newcommand{\class}{Math 629 Spring 26}
\newcommand{\term}{Fall 24}
\newcommand{\examnum}{Homework \hwnum}
\newcommand{\examdate}{\duedate}
\newcommand{\timelimit}{75 Minutes}
\newcommand{\vc}[3]{\langle #1,#2,#3\rangle}
\newcommand*{\vv}[1]{\vec{\mkern0mu#1}}
\newcommand{\bv}[1]{\boldsymbol{#1}}
\newcommand{\hide}[1]{}
\newcommand{\uvec}[1]{\boldsymbol{\hat{\textbf{#1}}}}
\newcommand{\vex}[1]{\boldsymbol{{\textbf{#1}}}}

\pagestyle{head}
\firstpageheader{}{}{}
\runningheader{\class}{\examnum\ - Page \thepage\ of \numpages}{\examdate}
\runningheadrule

\makeatletter
\renewcommand*\env@matrix[1][*\c@MaxMatrixCols c]{%
  \hskip -\arraycolsep
  \let\@ifnextchar\new@ifnextchar
  \array{#1}}
\makeatother

%\printanswers
\begin{document}

\noindent
\begin{tabular*}{\textwidth}{l @{\extracolsep{\fill}} r @{\extracolsep{6pt}} l}
\textbf{\class} & \textbf{Name: Mengxiang Jiang}\\
\end{tabular*}\\
\rule[2ex]{\textwidth}{2pt}

\newcommand{\duedate}{04/6/2026 at 11:59pm}
\newcommand\hwnum{10}

Turn this assignment in on Gradescope. All work must be neat and legible. Turn in 
a ``final draft''. There should be nothing scratched out and your work should flow
generally from left to right and top to bottom.

\begin{questions}


\question
Chapter 17 discusses Huygens's work on evolutes and involutes. Let the cycloid be
given parametrically by
\[
x=a(t-\sin t), \qquad y=a(1-\cos t).
\]
Show that
\[
\frac{dy}{dx}=\frac{\sin t}{1-\cos t}=\cot\!\left(\frac t2\right).
\]
Then compute the element of arclength $ds$ in terms of $dt$.
\\\\
First, we compute $dx$ and $dy$:
\[
  dx = a(1-\cos t)dt, \qquad dy = a\sin t dt
\]
So,
\[
  \frac{dy}{dx} = \frac{a\sin t dt}{a(1-\cos t)dt} = \frac{\sin t}{1-\cos t} = \cot\left(\frac t2\right).
\]
For the element of arclength, we have
\begin{align*}
  ds &= \sqrt{(dx)^2 + (dy)^2} \\
     &= \sqrt{(a(1-\cos t)dt)^2 + (a\sin t dt)^2} \\
     &= a\sqrt{(1-\cos t)^2 + \sin^2 t} dt \\
     &= a\sqrt{1 - 2\cos t + \cos^2 t + \sin^2 t} dt \\
     &= a\sqrt{2 - 2\cos t} dt \\
\end{align*}

\newpage 
\question
Using modern techniques, put Johann Bernoulli's solution
to the catenary problem,
$dy = \frac{a}{\sqrt{x^2 - a^2}}dx$
into a closed-form expression.
\\\\
The key identity that we will use is the following:
\[
  \frac{d}{dx} \cosh^{-1} x = \frac{1}{\sqrt{x^2 - 1}}.
\]
First, we can factor out $a$ from the denominator:
\[
  dy = \frac{a}{\sqrt{x^2 - a^2}} dx = \frac{a}{\sqrt{a^2(\frac{x^2}{a^2} - 1)}} dx = \frac{1}{\sqrt{\frac{x^2}{a^2} - 1}} dx.
\]
Now, we can make the substitution \( u = \frac{x}{a} \), which gives us \( du = \frac{1}{a} dx \) or \( dx = a du \). Substituting this into the integral, we get:
\[
  \int dy = \int \frac{1}{\sqrt{u^2 - 1}} a du = a \int \frac{1}{\sqrt{u^2 - 1}} du.
\]
Using the identity, we can integrate:
\[
  y = a \int \frac{1}{\sqrt{u^2 - 1}} du = a \cosh^{-1} u + C = a \cosh^{-1} \left( \frac{x}{a} \right) + C.
\]


\newpage 
\question 
Show why Lagrange's power series representation fails for
the case $f(x) =e^{-1/x^2}$.
\\\\
Suppose for contradiction that $f(x) = e^{-1/x^2}$ has a power series representation around $x = 0$. This means we can write:
\[
  f(x) = \sum_{n=0}^{\infty} a_n x^n,
\]
where $a_n$ are the coefficients of the power series. To find these coefficients, we can use the formula:
\[
  a_n = \frac{f^{(n)}(0)}{n!}.
\]
Now, we need to compute the derivatives of $f(x)$ at $x = 0$.
The first derivative is:
\[
  f'(x) = e^{-1/x^2} \cdot \frac{d}{dx}(-1/x^2) = e^{-1/x^2} \cdot \frac{2}{x^3}.
\]
In general the $n$-th derivative of $f(x)$ will be of the form:
\[
  f^{(n)}(x) = P_n\left(\frac{1}{x}\right) e^{-1/x^2},
\]
where $P_n(x)$ is some polynomial in $x$. As $x \to 0$, the term $e^{-1/x^2}$ approaches zero faster than any polynomial can grow. Therefore, we have:
\[
  \lim_{x \to 0} f^{(n)}(x) = 0.
\]
This means that all the coefficients $a_n$ in the power series representation will be zero:
\[
  a_n = \frac{f^{(n)}(0)}{n!} = 0.
\]
Thus, the power series representation of $f(x)$ around $x = 0$ is
\[
  f(x) = \sum_{n=0}^{\infty} 0 \cdot x^n = 0.
\]
However, this is a contradiction because $f(x) = e^{-1/x^2}$ is not identically zero for all $x$.
\end{questions}
\end{document}